\documentclass[11pt]{article}
\usepackage[margin=1in]{geometry}
\usepackage{amsmath,amssymb,amsthm}
\usepackage{graphicx}
\usepackage[hidelinks]{hyperref}

\newtheorem*{theorem}{Partial Theorem}
\newtheorem*{corollary}{Corollary}

\title{Partial Result: Complemented \(\ell_p\) Subspaces Give Singular Extensions}
\author{Automatic Math Discovery Run \texttt{fa\_banach\_001}}
\date{June 18, 2026}

\begin{document}
\maketitle

\section*{Source Problem}

The source paper is Spiros A. Argyros, Antonis Manoussakis, and Pavlos
Motakis, \emph{The HI extension of the standard HI spaces},
arXiv:2407.20030 \cite{AMM2024}. In its final discussion, the paper states
W. B. Johnson's singular-extension problem. In the parsed source this is
\path{data/parsed/arxiv_sources/2407.20030/source.tex}, lines 2026--2037.
On PDF page 25, Definition 10.1 defines a singular extension, and Problem 3
asks whether every infinite-dimensional separable Banach space not isomorphic
to \(c_0\) admits one.

\begin{center}
\includegraphics[width=0.92\linewidth]{figures/open_problem_crop.png}
\end{center}

\section*{Definitions}

Let \(X\) be a Banach space and let \(Y\) be a superspace of \(X\). Following
the source paper, \(Y\) is a \emph{singular extension} of \(X\) if \(Y/X\) is
infinite dimensional and the quotient map
\[
Q:Y\to Y/X
\]
is strictly singular. Equivalently, \(Q\) is not an isomorphism on any
infinite-dimensional subspace of \(Y\).

A subspace \(E\) of a Banach space \(X\) is complemented if
\(X=E\oplus F\) for some closed subspace \(F\), or equivalently if there is a
bounded projection from \(X\) onto \(E\).

\section*{Idea of the Proof}

Singular extensions are stable under adding an external complemented summand.
If \(E\) already has a singular extension \(0\to E\to Z\to V\to0\), then
\(X=E\oplus F\) is extended by \(Y=Z\oplus F\). The quotient still only sees
the \(Z\)-coordinate. Thus, any infinite-dimensional subspace on which the new
quotient map is bounded below would project injectively to an
infinite-dimensional subspace of \(Z\), where the old quotient map is also
bounded below. That contradicts the singularity of the original extension.

\section*{Complemented-Summand Permanence}

\begin{theorem}
Let \(X\) be a Banach space containing a complemented subspace \(E\). If
\(E\) admits a singular extension, then \(X\) admits a singular extension.
\end{theorem}

\begin{proof}
Choose a closed complement \(F\) so that \(X=E\oplus F\). Let
\[
0\longrightarrow E \xrightarrow{\iota} Z \xrightarrow{q} V
\longrightarrow 0
\]
be a singular extension of \(E\). Thus \(V\) is infinite dimensional,
\(\ker q=\iota(E)\), and \(q\) is strictly singular.

Let \(Y=Z\oplus_1 F\). Embed \(X=E\oplus F\) into \(Y\) by
\[
J(e+f)=(\iota e,f).
\]
Define
\[
\widetilde q:Y\to V,\qquad \widetilde q(z,f)=qz.
\]
Then \(\widetilde q\) is a bounded surjection and
\(\ker \widetilde q=J(X)\). Hence \(Y/J(X)\) is isomorphic to \(V\), so the
quotient is infinite dimensional.

It remains to show that \(\widetilde q\) is strictly singular. Suppose not.
Then there is an infinite-dimensional subspace \(M\subset Y\) and a constant
\(c>0\) such that
\[
\|\widetilde q y\|\ge c\|y\| \qquad (y\in M).
\]
Let \(P_Z:Y\to Z\) be the coordinate projection. If \(y\in M\) and
\(P_Zy=0\), then \(\widetilde qy=0\), hence \(y=0\). Thus \(P_Z|_M\) is
injective, and \(P_ZM\) is infinite dimensional. For \(z=P_Zy\in P_ZM\),
\[
\|qz\|=\|\widetilde qy\|\ge c\|y\|\ge c\|P_Zy\|=c\|z\|,
\]
where the inequality \(\|y\|\ge \|P_Zy\|\) uses the \(\ell_1\)-sum norm.
Therefore \(q\) is
bounded below on the infinite-dimensional subspace \(P_ZM\), contradicting
that \(q\) is strictly singular. Hence \(\widetilde q\) is strictly singular,
and \(Y\) is a singular extension of \(X\).
\end{proof}

\begin{corollary}
Let \(X\) be a separable Banach space. If \(X\) contains a complemented copy
of \(\ell_p\) for some \(1<p<\infty\), then \(X\) admits a separable singular
extension.
\end{corollary}

\begin{proof}
Kalton and Peck proved that \(\ell_p\), \(1<p<\infty\), admits a singular
extension \(Z_p\) \cite[Theorem 6.4]{KP1979}; the source paper recalls this
example immediately after Johnson's Problem 3. Transport this extension across
the isomorphism between the complemented subspace of \(X\) and \(\ell_p\), and
apply the theorem. Separability is preserved because \(X\), the complement,
and the Kalton--Peck extension \(Z_p\) are separable.
\end{proof}

\section*{Scope and Limitations}

This gives a positive answer to Johnson's Problem 3 for the concrete class of
separable Banach spaces containing a complemented copy of \(\ell_p\),
\(1<p<\infty\). It is not a solution of Johnson's full problem, since it does
not cover arbitrary separable spaces not isomorphic to \(c_0\). It is also not
a solution of Pelczynski's HI-extension problem, because HI spaces do not
decompose into nontrivial complemented infinite-dimensional summands.

\section*{Novelty and Search Bounds}

The run indexes were searched for arXiv:2407.20030, HI extension, singular
extension, Kalton--Peck, \(Z_p\), pushout, pullback, twisted sum, and strictly
singular quotient-map terminology. No prior packet for this source paper or
for this complemented-summand permanence lemma was found. The result is a
small but useful permanence observation, not a claimed new construction of
singular extensions beyond the complemented-subspace mechanism.

\section*{Human Review Recommendation}

Check the permanence lemma against the exact convention for singular twisted
sums. The proof uses the quotient-map definition from the source paper and the
\(\ell_1\)-direct sum only to make the projection estimate immediate.

\begin{thebibliography}{9}
\bibitem{AMM2024}
S. A. Argyros, A. Manoussakis, and P. Motakis,
\emph{The HI extension of the standard HI spaces}, arXiv:2407.20030, 2024.

\bibitem{KP1979}
N. J. Kalton and N. T. Peck,
\emph{Twisted sums of sequence spaces and the three space problem},
Trans. Amer. Math. Soc. 255 (1979), 1--30.
\end{thebibliography}

\end{document}
